% Verbatim copy of /Users/alingalatan/workspaces/grothe/handoff/manuscript/legacy/ch34-gpt-shared-nongaussian.tex (file dated 2026-10-25). Only this line was added.
% Consolidated source module (formerly Chapter 34). Owner: lower-bound/10-coupled-complements.tex
\section{Lower-bound GPT: shared non-Gaussian $A/B$ couplings and a quantitative weak form of K (20 October)}
\label{ch:shared-nongaussian-new}
\chapsummary{Shared non-Gaussian comparisons, independent-block bounds and joint replacement tools. Smaller regions remain dependencies; current coverage is in \autoref{ch:K-state}.}

\paragraph{Provenance and status (IDE note, 20 October).} Lower-bound GPT's \code{proof/ch34-shared-nongaussian.tex} from \pathref{gpt-reviews/shared-nongaussian-20261020/}, ported
verbatim apart from this header. \emph{Our replay} \st{C}: \code{python3 -I verify\_all.py} from an unrelated directory: \textbf{PASS complete portable replay in 152.5\,s}, including
the inherited Gaussian-covariance certificate and the characteristic-function inequality. The signed Hankel signs are the visible-averaging principle of \autoref{ch:forest} in
Bessel form. Item (iii) is a block version of the replacement gap $D$ of \autoref{ch:polarized-profile-certificate}.

\noindent At the peak, write $\Lambda(R)=\mathbb E(|\ell(T)+R|+|\ell(T)-R|)-\mathbb E|R|^2/\varepsilon$, and $S=\sup_{v\ge0}\{2\mathbb E|\ell(T)+\sqrt vG|-v/\varepsilon\}$. All new results below concern this exact profile. The independent benchmark certificate gives $S>1.37395328457$.\par


\subsection{Finite shared B and simultaneous A/B couplings}
\label{sec:finite-shared-ab}
This section proves new finite-amplitude, non-Gaussian cases of
Conjecture K. It does not change the global Grothendieck endpoint.
Keep the exact peak parameters
\[
 a=699/1250,\quad b=179/1250,\quad \varepsilon=1891/2500,
 \quad \ell(t)=\sqrt t(a+bt),\quad A(t)=t-1,\quad B(t)=t/\sqrt2.
\]
Write $H(X)=X(|X|^2-2)/\sqrt2$, and let $S$ be the diffuse Gaussian
benchmark. All Gaussians below are independent standard proper complex
Gaussians unless a shared driver is explicitly displayed.

\begin{theorem}[Unequal Wick sums with shared B]
Let $q=\sum_jc_jH(X_j)$ and
$R=q+\sqrt vG+B(T)\sum_j\beta_j\overline{X_j}$, where
$T=|Z|^2$. Put $u=\sum|c_j|^2$, $c_*=\max|c_j|$,
$B_4=\sum|c_j|^4$, and $r=\|\beta\|_2$. If
\[
 7/200\le v\le9/200,\qquad u\le1/50,\qquad
 c_*\le1/25,\qquad r\le1/100,
\]
then
\[
 \Lambda(R)\le S-\frac{27}{10}B_4-\frac{r^2}{20}.
\]
The number of coordinates and all coefficient phases are unrestricted.
\end{theorem}

\begin{theorem}[Simultaneous shared A/B]
Let
\[
 R=\sum_jc_jH(X_j)+\sqrt vG
   +A(T)\sum_j\alpha_jX_j+B(T)\sum_j\beta_j\overline{X_j},
 \qquad r^2=\|\alpha\|_2^2+\|\beta\|_2^2.
\]
If
\[
 7/200\le v\le9/200,\qquad
 \sum_j|c_j|^2\le16/2500,\qquad
 \max_j|c_j|\le1/50,\qquad r\le1/1000,
\]
then
\[
 \boxed{\Lambda(R)\le S-\frac34\sum_j|c_j|^4-\frac{r^2}{200}.}
\]
There is no restriction on the relative directions of the two driver
coefficient vectors.
\end{theorem}
\paragraph{Update (IDE note, 20 October evening).} Theorem~\ref{thm:buffered-review-anyn} (\autoref{ch:buffered-drivers}) removes the restriction
$r\le1/1000$ on the same variance, coefficient and buffer range and raises the quartic coefficient to $2$ (and allows $e\le1/100$ with $49/100$), at the price of the
weaker $A/B$ penalty $E_{AB}/2000$. On $r\le1/1000$ the penalty $r^2/200$ above remains the stronger statement. The visible-kernel signs of this chapter are inputs
to the new theorem.

\paragraph{Scope.}
The B theorem allows twelve coefficients $c_j=1/25$ and $v=7/200$;
their normalized fourth moment is $179138/73441>2.439$.
The simultaneous theorem allows sixteen coefficients $c_j=1/50$ and
$v=7/200$; their normalized fourth moment exceeds $2.0627$.
Thus both cover genuine shared cubic examples in the application's
kurtosis range. The Gaussian buffer and finite-amplitude restrictions
remain essential to the present proof.

\subsubsection{Positive Hankel representations after visible averaging}
Put $h(t)=t(a+bt)^2$, $d(t)=h'(t)$, and $\tau=h^{-1}$.
The functions
\[
 b_1(h(t))=\frac{e^{-t}}{(a+bt)^2},\quad
 b_2(h(t))=\frac{e^{-t}}{\sqrt2(a+bt)^3(a+3bt)},\quad
 b_4(h(t))=\frac{e^{-t}}{2(a+bt)^5(a+3bt)}
\]
are completely monotone. Indeed $d^{-1}$ is completely monotone in
$t$, and the operation $-d^{-1}\partial_t$ preserves that property.
Induction proves all alternating derivative inequalities after
composition with $\tau$. Write $b_j(y)=\int e^{-sy}\mu_j(ds)$ with
positive measures.

For
$\Psi_A(k)=\mathbb E[A(T)J_0(k\ell(T))]$ and
$\Psi_B(k)=\mathbb E[B(T)J_2(k\ell(T))]$, integration by parts and the
elementary Gaussian Hankel integrals give
\[
 \Psi_A(k)=-\frac{k^2}{4}\int s^{-2}e^{-k^2/(4s)}\mu_1(ds)\le0,
 \qquad
 \Psi_B(k)=\frac{k^2}{4}\int s^{-3}e^{-k^2/(4s)}\mu_2(ds)\ge0.
\]
In particular,
\begin{align*}
 I_A:=\int_0^\infty k^2|\Psi_A(k)|\,dk
 &=3\int_0^\infty
 \frac{e^{-t}(a+bt+2b)}{\sqrt t(a+bt)^4}\,dt,\\
 I_B:=\int_0^\infty k^2\Psi_B(k)\,dk
 &=\frac3{\sqrt2}\int_0^\infty
 \frac{e^{-t}}{\sqrt t(a+bt)^3}\,dt<17.
\end{align*}
Let $R_X=|X|^2$ and $h_X=\sqrt{R_X}(R_X-2)/\sqrt2$. The content
integrals contain $\mathbb E[\sqrt{R_X}J_1(kch_X)]/k$.
Use
\[
 \mathbb E\sqrt{R_X}h_X=0,\quad
 |J_1(x)-x/2|\le|x|^3/16,\quad
 \mathbb E\sqrt{R_X}|h_X|^3
 =\frac{-32+480e^{-2}}{2\sqrt2}.
\]
The exact scalar checker proves
$I_A\mathbb E\sqrt{R_X}|h_X|^3/8<51$ and
$I_B\mathbb E\sqrt{R_X}|h_X|^3/8<25$. Consequently, with
$B_6=\sum|c_j|^6$,
\[
 \|\nabla_\alpha\Lambda|_0\|_2\le51\sqrt{B_6},\qquad
 \|\nabla_\beta\Lambda|_0\|_2\le25\sqrt{B_6}.
\]
These bounds hold with any independent circular buffer: its radial
characteristic function has absolute value at most one, and the other
Wick blocks can be included in that buffer.

\subsubsection{B curvature and the mixed entry}
For $\psi_g(x)=\mathbb E|x+gG|$, define
\[
 M_B(x)[p,p]=\mathbb E_{T,\theta}B(T)^2
 D^2\psi_g(\ell(T)e^{i\theta}+x)
 [e^{2i\theta}\bar p,e^{2i\theta}\bar p].
\]
This field is even in $x$. The harmonic-zero and harmonic-four
transforms of its visible source are, with $z=k^2/(4s)$,
\[
 \Psi_0(k)=\int s^{-3}e^{-z}(2-4z+z^2)\mu_4(ds),\quad
 \Psi_4(k)=\int s^{-3}e^{-z}z^2\mu_4(ds)\ge0.
\]
Fourier differentiation and the angular integral give
\[
 |M_B(x)[p,p]-M_B(0)[p,p]|
 \le\frac{|x|^2|p|^2}{8}
       \int_0^\infty k^2(|\Psi_0|+\Psi_4)\,dk
 <10|x|^2|p|^2.
\]
For the last inequality use
$|2-4z+z^2|+z^2\le2+4z+2z^2$ and
\[
 \frac18\int k^2(|\Psi_0|+\Psi_4)\,dk
 \le\frac{31}{8}\int_0^\infty
        \frac{\sqrt t e^{-t}}{(a+bt)^3}\,dt
 <\frac{31}{8}\frac52<10.
\]
The Gaussian convolution preserves this uniform derivative bound.

The analogous mixed field has weight $AB$ and one direction $p$,
the other $e^{2i\theta}\bar z$. Its source is $z^2b_{AB}(|z|^2)/\pi$,
where
\[
 b_{AB}(h(t))=
 \frac{(t-1)e^{-t}}{\sqrt2(a+bt)^3(a+3bt)}.
\]
It is the difference $b_{(2)}/b-(a/b+1)b_{(3)}$ of two completely
monotone functions, where
$b_{(p)}(h(t))=e^{-t}/[\sqrt2(a+bt)^p(a+3bt)]$.
Set
\[
 I_p=\frac3{\sqrt2}\int_0^\infty
          \frac{e^{-t}}{\sqrt t(a+bt)^p}\,dt.
\]
The angular estimate
$\int\cos^2(\phi-e)|\cos(\phi-p)\cos(\phi-z)|\,d\phi\le3\pi/4$
and evenness of the mixed field yield
\[
 |M_{AB}(x)-M_{AB}(0)|
 \le\frac3{16}\left[I_2/b+(a/b+1)I_3\right]|x|^2|p||z|
 <29|x|^2|p||z|.
\]
This is a direct mixed estimate; it is not inferred from diagonal
bounds for a possibly indefinite Hessian difference.

\subsubsection{Finite Taylor estimates and completion}
Put $D=u+3c_*^2$. The exact Wick identity
$\mathbb E|q|^2|P|^2\le D\|P\|_2^2$ holds for every proper linear
driver $P$. The earlier averaged A curvature costs at most $12D$;
the new B and mixed costs are $10D$ and $29D$.
The Gaussian local certificates give diagonal penalties $2/5,9/10$
and signed mixed coefficient of absolute value at most $17/100$.
At $D\le19/2500$, the matrix
\[
 \begin{pmatrix}2/5-12D&-(17/100+29D)\\
 -(17/100+29D)&9/10-10D\end{pmatrix}
 \succeq\frac{49}{500}I.
\]
All principal-minor comparisons are rational.

The inherited derivative bound is $\|D^4\psi_g\|\le204$.
For a normalized pair of proper linear drivers, put
$P_T=AP+B\bar Q$, $\|P\|_2^2+\|Q\|_2^2=1$.
Circle charge and Gaussian moments give
\[
 \mathbb E|q||P_T|^3\le\frac52\sqrt{3D},\qquad
 \mathbb E|P_T|^4\le27.
\]
The odd third derivative vanishes at $q=0$. Taylor's theorem therefore
proves
\[
 \Lambda(R)\le\Lambda(q+\sqrt vG)+\sqrt{3226B_6}\,r
 -\left(\frac{49}{500}-170\sqrt{3D}\,r-459r^2\right)r^2.
\]
For $r\le1/1000$ the bracket exceeds $1/200+1/15$.
The earlier unequal-Wick Fourier deficit gives
$\Lambda(q+\sqrt vG)\le S-(17/3)B_4$ in this parameter range.
Completing the square costs at most
$(15/4)3226B_6\le4.839B_4$, leaving more than $3B_4/4$ and $r^2/200$.

For B alone, the sharper remainder is
\[
 \Lambda(R)\le\Lambda(q+\sqrt vG)+25\sqrt{B_6}\,r
 -\left(\frac9{10}-10D-204\sqrt D\,r-204r^2\right)r^2.
\]
The fresh elementary estimate
$\phi_H(s)\ge\exp(-s^2/4+(2/5)s^4)$ for $0\le s\le1/5$,
together with the inherited $\phi_H\ge e^{-s^2/4}$ on $[0,4]$,
gives a pure deficit at least $(37/10)B_4$ for $u\le1/50,c_*\le1/25$.
Its low-frequency integral is bounded below by
\[
 \int_0^5 k^2e^{-(33/160)k^2}\,dk>93/20;
\]
the complete high-frequency loss is below $B_4/1000$.
On the B theorem's box its curvature bracket exceeds $3/10$.
Reserve $r^2/20$ and use
$25\sqrt{B_6}r-r^2/4\le625B_6\le B_4$.
This proves the first theorem.

\paragraph{Verification and dependencies.}
The standalone standard-library script \texttt{check\_b\_constants.py}
checks every new scalar bound using rational upper sums, rational
exponential bounds, and exact polynomial comparisons. The full
derivations are in \texttt{B\_SECTOR\_PROOF.md} and
\texttt{JOINT\_AB\_PROOF.md}. The inherited characteristic-function,
A-curvature, and Gaussian-derivative certificates are explicit inputs.
The signed Gaussian mixed-coefficient bound is supplied by this round's
\texttt{polarized-audit-20261020/check\_support.cpp}.

\subsubsection{Larger B-only coefficient ranges}
Retaining the absolute Laguerre integral sharpens the B curvature
constant from $10$ to $17/5$. The exact positive upper sums prove
\[
 \int_0^\infty \sqrt z e^{-z}|z^2-4z+2|\,dz<3/2,
 \qquad \int s^{-3/2}\mu_4(ds)<7/5.
\]
Since $\Gamma(7/2)<10/3$, the preceding Fourier bound is at most
$\tfrac12(7/5)(3/2+10/3)<17/5$. Hence globally
\[
 \Lambda(R)\le\Lambda(q+\sqrt vG)+25\sqrt{B_6}r
 -\left(\frac9{10}-\frac{17}{5}D-204\sqrt D\,r-204r^2\right)r^2.
\]
For $u\le1/50$ and $v\in[7/200,9/200]$ this yields
\begin{align*}
 c_*\le3/50,\ r\le1/100
 &\quad\Longrightarrow\quad
 \Lambda(R)\le S-3B_4/2-r^2/40,\\
 c_*\le2/25,\ r\le1/1000
 &\quad\Longrightarrow\quad
 \Lambda(R)\le S-B_4/2-r^2/100.
\end{align*}
The pure-deficit low bands are respectively $[0,10/3]$ and $[0,5/2]$.
Exact Taylor lower integration gives
\[
 \int_0^{10/3}k^2e^{-(33/160)k^2}\,dk>94/25,\qquad
 \int_0^{5/2}k^2e^{-(33/160)k^2}\,dk>127/50.
\]
After the separately checked infinite-frequency allowance $B_4/1000$,
the pure deficits exceed $3B_4$ and $2B_4$. The finite-coupling
curvatures exceed $1/40+3/8$ and $18/25$. Completing squares costs
at most $3B_4/2$ and $100B_4/71$, respectively.
Five coefficients $3/50$ with $v=7/200$ have normalized fourth moment
about $2.969$; three coefficients $2/25$ give about $3.75684$.
Both are genuine finite shared-B cases in the desired kurtosis range.
They do not enlarge the simultaneous A/B box without further work.


\subsection{A fourth-order replacement theorem for independent shared-driver blocks}
\label{sec:gamma-block-replacement}

This section proves a sufficient form of Conjecture K for independent
blocks whose cubics and linear coupling drivers may have arbitrary
internal dependence.  The replacement error is fourth order in each
block, so equal blocks cost $N^{-1}$ rather than $N^{-1/2}$.
A concrete shared A/B family lies in the application's kurtosis window.
This does not prove K for general coupled kernels and does not change
the Grothendieck endpoint by itself.  The numerical Gaussian covariance
theorem proved in the preceding polarized-profile package is an input.

\subsubsection{An exact scalar identity}
Let $S$ be an odd polynomial of finitely many independent standard real
Gaussians. Write $\mathcal L$ for their Ornstein--Uhlenbeck generator and
$\mathcal A=-\mathcal L^{-1}$, with constants sent to zero. Put
\[
 v=\mathbb E S^2,\qquad
 \tau=\langle DS,D\mathcal AS\rangle,\qquad H=\tau-v,
\]
\[
 K=\langle DS,D\mathcal AH\rangle,\qquad
 J=\langle DS,D\mathcal AK\rangle.
\]
Here $H$ is centered and even, while $K$ is odd and hence centered.
If $N$ is an independent real Gaussian of variance $v$, then
\begin{equation}\label{eq:block-interpolation}
 \mathbb Ef(S)-\mathbb Ef(N)
 =\frac12\int_0^1 s\,
  \mathbb E\left[Jf^{(4)}(\sqrt{s}S+\sqrt{1-s}N)\right]ds
\end{equation}
for every $C^4$ test with bounded fourth derivative and polynomially
bounded lower derivatives. In particular,
\[
 |\mathbb Ef(S)-\mathbb Ef(N)|
 \le\frac14\|f^{(4)}\|_\infty\mathbb E|J|.
\]
To prove the identity, set $Y_s=\sqrt{s}S+\sqrt{1-s}N$.
Coordinate Gaussian integration by parts gives
\[
 \frac d{ds}\mathbb Ef(Y_s)=\frac12\mathbb E[Hf''(Y_s)].
\]
For any centered polynomial $P$, Hermite expansion and coordinate
integration by parts give
$\mathbb E[PF]=\mathbb E\langle D\mathcal AP,DF\rangle$.
Applying this first to $H$, then to $K$, yields
\[
 \mathbb E[Hf''(Y_s)]=\sqrt{s}\mathbb E[Kf'''(Y_s)]
                   =s\mathbb E[Jf^{(4)}(Y_s)].
\]
Integration proves~\eqref{eq:block-interpolation}; approximation handles
the endpoints. Degenerate $v=0$ is trivial. The normalization is
\[
 \mathbb EJ=\frac{\mathbb ES^4-3v^2}{6},
\]
since $\mathbb ES^4=3\mathbb E[\tau S^2]$,
$\mathbb E[HS^2]=2\mathbb E[KS]$, and $\mathbb E[KS]=\mathbb EJ$.
The polynomial $J$ need not be nonnegative, so its absolute expectation
cannot simply be replaced by its mean. This is the explicit
iterated-Gamma mechanism in smooth cumulant estimates; compare
Bierm\'e--Bonami--Nourdin--Peccati, arXiv:1109.1546.

\subsubsection{A joint block comparison, with the interpolating Gaussian retained}
Use
\[
 \epsilon=\frac{1891}{2500},\quad
 \ell(t)=\sqrt t\left(\frac{699}{1250}+\frac{179}{1250}t\right),
 \quad A(t)=t-1,\quad B(t)=t/\sqrt2,
\]
and the Gaussian benchmark
\[
 S_* =\sup_{v\ge0}\left\{2\mathbb E|\ell(T)+\sqrt vG|-v/\epsilon\right\},
 \qquad T\sim\operatorname{Exp}(1).
\]
On mutually independent Gaussian blocks take
$q_j\in H_{2,1}$, $U_j\in H_{1,0}$, $V_j\in H_{0,1}$.
All internal dependence is permitted. Set
\[
 R_T=gG+\sum_j(q_j+A(T)U_j+B(T)V_j),\qquad g>0,
\]
\[
 p=\sum_j\|q_j\|_2^2,\quad
 \alpha^2=\sum_j\|U_j\|_2^2,\quad
 \beta^2=\sum_j\|V_j\|_2^2.
\]
For each $t,\theta$ construct $J_{j,t,\theta}$ as above from
\[
 S_{j,t,\theta}=\operatorname{Re}\left(e^{-i\theta}
                   (q_j+A(t)U_j+B(t)V_j)\right).
\]
Angular brackets denote uniform averaging over $[0,2\pi]$. Write
\[
 \mathcal B_j=\mathbb E_T\langle\mathbb E|J_{j,T,\theta}|\rangle,
 \qquad
 C(g,p)=\frac{4\sqrt\pi}{g(g+\sqrt{g^2+p})^2}.
\]
If $\mathcal D$ is the signed difference between the two-sign norm
payoff and its matched Gaussian payoff, then
\begin{equation}\label{eq:block-error}
 \boxed{\mathcal D\le C(g,p)\sum_j\mathcal B_j.}
\end{equation}
The Gaussian covariance theorem therefore implies
\begin{equation}\label{eq:block-K}
 \boxed{\Lambda(R)\le S_*-\alpha^2/10-\beta^2/4
                         +C(g,p)\sum_j\mathcal B_j.}
\end{equation}

Indeed, the quantities $H,K,J$ add over independent blocks: mixed
coordinate gradients vanish and the inverse generator is additive.
The total scalar variance $v_\theta(t)$ is at least $p/2$, since the
cubic part is proper and orthogonal to its linear drivers. At
interpolation time $s$, retain the independent interpolating Gaussian
inside the smoothing variance, which is
\[
 g^2/2+(1-s)v_\theta(t)\ge[g^2+(1-s)p]/2.
\]
For a real standard Gaussian $N_0$,
\[
 \left\|\frac{d^4}{dx^4}\mathbb E|x+\sigma N_0|\right\|_\infty
 =\sqrt{2/\pi}\,\sigma^{-3}.
\]
This follows directly from
$\sup_u|(u^2-1)e^{-u^2/2}|=1$. The projection identity
$2|z|=\pi\langle|\operatorname{Re}(e^{-i\theta}z)|\rangle$,
odd symmetry, and~\eqref{eq:block-interpolation} give
\[
 \mathcal D\le
 2\sqrt\pi\int_0^1\frac{s\,ds}{[g^2+(1-s)p]^{3/2}}
 \sum_j\mathcal B_j.
\]
The integral equals $2/[g(g+\sqrt{g^2+p})^2]$, also at $p=0$ by
continuity. This proves~\eqref{eq:block-error}. Keeping this Gaussian
improves the cruder coefficient $\sqrt\pi/g^3$ by the factor
$4g^2/(g+\sqrt{g^2+p})^2$.

Chaos orthogonality gives the matched Gaussian covariance
\[
 V(t)=g^2+p+\alpha^2A(t)^2+\beta^2B(t)^2,
 \qquad |P(t)|\le2\alpha\beta|A(t)B(t)|.
\]
Thus it belongs to the proved Gaussian covariance class; its price is
unchanged, proving~\eqref{eq:block-K}. For $N$ equal blocks scaled by
$N^{-1/2}$, the last polynomial scales by $N^{-2}$, so its total
budget is the one-block budget divided by $N$.

\subsubsection{An explicit Wick polynomial and its budget}
Let $X=(x+iy)/\sqrt2$ and
$H(X)=X(|X|^2-2)/\sqrt2$, where $x,y$ are standard real Gaussians.
For the scalar polynomial
\[
 S=c\operatorname{Re}H(X)+ax+by,
 \qquad \operatorname{Re}H(X)=(x^3+xy^2-4x)/4,
\]
the last polynomial in~\eqref{eq:block-interpolation} is exactly
\[
 J=ca(a^2+b^2)+c^2Q_{22}+c^3(aP+bQ)+c^4R,
\]
where
\begin{align*}
Q_{22}={}&(27a^2/16+61b^2/144)x^2
 +(61a^2/144+3b^2/16)y^2+(37ab/36)xy\\
 &+7a^2/18-b^2/9,\\
P={}&27x^4/32+29x^2y^2/45+27x^2/40
       +73y^4/1440+91y^2/360-23/90,\\
Q={}&xy(301x^2+121y^2-56)/720,\\
R={}&27x^6/256+577x^4y^2/3840+23x^4/80
 +179x^2y^4/3840+7x^2y^2/40+67x^2/240\\
 &+7y^6/3840+y^4/80+49y^2/240-17/15.
\end{align*}
Exact finite polynomial expansion verifies this identity and gives
\[
 \mathbb E|Q_{22}|\le\frac52a^2+\frac{13}{18}b^2,
 \quad \mathbb E|P|\le\frac{203}{45},
 \quad \mathbb E|Q|\le\frac5{2\pi},
 \quad \mathbb E|R|\le\frac{587}{120}.
\]
For $Q_{22}$, its quadratic part is positive semidefinite: its
$144^2$-scaled determinant is
$14823a^4+4806a^2b^2+1647b^4$.
Its constant is at least $-b^2/9$ and its expectation is
$5a^2/2+b^2/2$. For $P,R$, every nonconstant coefficient is
nonnegative; their means are $4,21/8$ and their constants are
$-23/90,-17/15$. Finally the absolute first and third Gaussian moments
prove the estimate for $Q$.

For the complex block
\[
 cH(X)+\sqrt2a_0A(T)X+2b_0B(T)\overline X,
 \qquad c\ge0,\quad a_0,b_0\in\mathbb C,
\]
rotation gives the scalar coefficients
\[
 a=u(T-1)+|b_0|T\cos2\theta,\qquad b=w(T-1)+|b_0|T\sin2\theta,
 \quad u^2+w^2=|a_0|^2.
\]
In the following budget the arguments $a_0,b_0$ denote their nonnegative
moduli. A valid one-block budget, uniform in all relative phases, is
\begin{align*}
 \mathfrak B(c,a_0,b_0)={}&c\left[
 \frac{12}{\pi}b_0^3+\frac{21}{\pi}a_0^2b_0
                          +(12/e-2)|a_0|^3\right]\\
 &+c^2\left[\frac52a_0^2+\frac{29}{9}b_0^2\right]\\
 &+c^3\left[\frac{46}{5e}|a_0|
       +\left(\frac{406}{45\pi}+\frac5{\pi^2}\right)b_0\right]
 +\frac{587}{120}c^4.
\end{align*}
The only nontrivial angular bound used here is
\[
 \left\langle|a+b\cos\theta|(a^2+2ab\cos\theta+b^2)\right\rangle
 \le|a|^3+\frac7\pi a^2b+\frac2\pi b^3,\qquad b\ge0.
\]
For $0\le r=|a|/b\le1$, the left side is $2b^3F(r)/\pi$, where
\[
 F(r)=(1+2r^2)\sqrt{1-r^2}+r(r^2+2)\arcsin r.
\]
Now $F(0)=1,F'(0)=0,F''(0)=7$ and
$F'''(r)=6\arcsin r-r/\sqrt{1-r^2}\le3\pi$; Taylor's integral
remainder proves the bound. For $r\ge1$ the exact answer is
$|a|^3+2|a|b^2$, and $\pi<7/2$ suffices. For a complex centre $x+iy$ of modulus $r$, the angular average of
$|\operatorname{Re}(x+iy+b e^{i\theta})|\,|x+iy+b e^{i\theta}|^2$
is $(r^2+b^2)I_0(x,b)+2bxI_1(x,b)$, where
$I_0=\langle|x+b\cos\theta|\rangle$ and
$I_1=\langle\cos\theta|x+b\cos\theta|\rangle$.
Both $I_0$ and $xI_1$ increase with $|x|$, so this is maximized at
$x=r,y=0$. For $0\le x\le b$,
$I_1=[x\sqrt{1-x^2/b^2}+b\arcsin(x/b)]/\pi$
and $I_1'=2\sqrt{1-x^2/b^2}/\pi\ge0$; for $x\ge b$, $I_1=b/2$.
The cubic A coefficient follows from
$\sqrt{(203/45)^2+(5/(2\pi))^2}<23/5$.
Thus the bound includes all relative phases. The required exponential
moments are
\[
 \mathbb E|T-1|=2/e,\quad \mathbb E|T-1|^3=12/e-2,
 \quad \mathbb E[(T-1)^2T]=3,\quad\mathbb ET^3=6.
\]

\subsubsection{Finite shared A/B examples in the application window}
Take independent proper Gaussians and set
\[
 q_N=\frac3{25\sqrt N}\sum_{j=1}^N H(X_j)+\frac9{50}G,
 \quad
 U_N=\frac{\sqrt2}{100\sqrt N}\sum_{j=1}^N X_j,
 \quad
 V_N=\frac{48}{125\sqrt N}\sum_{j=1}^N\overline{X_j}.
\]
Arbitrary unit phases may multiply the A and B coefficients, independently
for each summand. The checker
proves exactly
\[
 C(9/50,(3/25)^2)\mathfrak B(3/25,1/100,24/125)
 <\frac{8987}{5000}.
\]
Hence
\[
 \boxed{\Lambda(q_N+AU_N+BV_N)
 \le S_*-\frac{9221}{250000}+\frac{8987}{5000N}
 <S_*-10^{-4}\quad(N\ge49).}
\]
The remote variance is $117/2500$, and its normalized fourth moment is
\[
 2+\frac{42(3/25)^4}{N(117/2500)^2}.
\]
It belongs to the stated window $[2.06,6]$ for $49\le N\le66$.
The B amplitude $.384$ provides a substantial profile deficit;
this does not settle the small-coupling obstruction near a maximizer.

There is also an explicit finite pure-cubic version. Replace the buffer
by $(9/(50\sqrt M))\sum_{k=1}^M H(Y_k)$, using independent $Y_k$.
The exact Wick Stein identity is
\[
 \mathbb E\|\Gamma_H-I_2/2\|_{\rm HS}^2=13/2,
\]
so the multivariate Gaussian Stein bound gives Wasserstein error at most
$\sqrt{13/M}$ for the normalized sum. The two-sign payoff is
$2$-Lipschitz and the variance price does not change. Exact squaring
verifies
\[
 2(9/50)\sqrt{13/20000}
 <\frac{9221}{250000}-\frac{8987}{5000\cdot66}-10^{-4}.
\]
Thus, writing this new pure-cubic remote residual as
$\widetilde q_{N,M}$,
\[
 \boxed{\Lambda(\widetilde q_{N,M}+AU_N+BV_N)<S_*-10^{-4}
                  \quad(N\ge66,\ M\ge20000).}
\]
At $N=66,M=20000$ its remote normalized fourth moment is
$2.0612539537\ldots$; for $N=66$ it remains in $[2.06,6]$ for all
larger $M$. The checker retains the larger rational deficit
$283/600000$: the buffered deficit is $3619/375000$ and the buffer
error is below $9179/1000000$. No high-dimensional numerical integration is used.
An optional calibration replaces the coefficient $48/125$ of $V_N$
by $3/5$. Its cost is below $243281/(50000N)$ and its profile penalty
is $4501/50000$; $N\ge66,M\ge7000$ then retain the same $10^{-4}$
deficit. Its normalized fourth moment at those minimal dimensions is
$2.0631231845\ldots$. Its retained rational deficit is
$5171/6600000$, subtracting buffer error $3103/200000$ from the
buffered deficit $10757/660000$.

\subsubsection{Verification and limitation}
The standard-library script \texttt{check\_gamma\_replacement.py}
reconstructs $J$ using exact rational Hermite expansions, verifies all
polynomial identities, positivity claims, moments and rational
constants, and passes under isolated Python mode from another
directory. It does not certify general coupled kernels, even Y/W
interactions, the diagonal term, or a full row envelope. The remaining
mathematical opportunity is to retain the signed pairing with $J$ in
\eqref{eq:block-interpolation} instead of replacing it by
$\mathbb E|J|$. The global Grothendieck endpoint remains unchanged.


\subsection{Structural consequences of the signed visible average}
\label{sec:signed-structural-ab}
The two preceding calculations distinguish two useful mechanisms.
For small shared couplings, the true first-order content is paid by the
non-Gaussian residual's deficit. For independent blocks, the cubic and
its linear drivers are replaced jointly, with a fourth-order error.
Neither mechanism uses Gaussian domination at each fixed visible shift;
that stronger statement is false.

\subsubsection{A completely monotone formula for the actual A profile}
The following identity holds for every $a>0,b\ge0$, not only the peak.
Let $T\sim\mathrm{Exp}(1)$, let $\theta$ be independent and uniform,
and put
\[
 h(t)=t(a+bt)^2,\qquad L=\sqrt{h(T)}e^{i\theta},\qquad
 F_A(u)=\E[(T-1)|L+x|],\quad u=|x|^2.
\]
As above, $b_1(h(t))=e^{-t}/(a+bt)^2$ is completely monotone;
write $b_1(y)=\int e^{-sy}\mu_1(ds)$ with $\mu_1\ge0$.
Then
\begin{equation}\label{eq:cm-actual-a-profile}
 \boxed{F_A(u)=\frac{\sqrt\pi}{4}\int s^{-3/2}
       e^{-su/2}I_0(su/2)\,\mu_1(ds).}
\end{equation}
In particular $F_A$ is positive, decreasing and completely monotone
in $u$, and
\[
 F_A(0)=\sqrt\pi(a/4+9b/8),\qquad F_A(u)\longrightarrow0.
\]
Indeed exponential integration by parts and the Gaussian Hankel integral give
\[
 \Psi_A(k):=\E[(T-1)J_0(k\sqrt{h(T)})]
 =-\frac{k^2}{4}\int s^{-2}e^{-k^2/(4s)}\mu_1(ds)\le0.
\]
Apply the planar norm Fourier identity to $F_A(u)-F_A(0)$ and use
\[
 \int_0^\infty e^{-k^2/(4s)}J_0(k\sqrt u)\,dk
   =\sqrt{\pi s}\,e^{-su/2}I_0(su/2).
\]
The integration constant is fixed by the limit at infinity.
For that limit, subtract $|x|$, use $\E(T-1)=0$, and dominate by
$|T-1||L|$; the pointwise directional limit has zero angular mean.
All required moments are finite. Finally
\[
 e^{-su/2}I_0(su/2)=\frac1{2\pi}\int_0^{2\pi}
                   e^{-su\cos^2\vartheta}\,d\vartheta
\]
exhibits a positive Laplace mixture and proves complete monotonicity.
The finite value at zero justifies the positive-mixture exchanges.
The standard integral identities are available in
\href{https://dlmf.nist.gov/10.22.E51}{DLMF 10.22.51} and
\href{https://dlmf.nist.gov/10.32.E1}{DLMF 10.32.1}.

This identity concerns the exact averaged A profile. It does not by
itself bound its first-chaos projection when $q$ and the A driver share
coordinates. In particular it must not be turned into a pointwise
Gaussian comparison for the whole perturbed payoff.

\subsubsection{Gaussian-buffer closure does not add a separate conjecture}
Let $\rho=\sqrt\pi/2$. The following two universal assertions are equivalent:
\begin{enumerate}
\item Every finite-dimensional $q\in H_{2,1}$ satisfies
      $\E|q|\le\rho\sqrt{\E|q|^2}$.
\item For every such $q$, every $v\ge0$ and an independent proper
      standard Gaussian $G$,
      $\E|q+\sqrt vG|\le\rho\sqrt{\E|q|^2+v}$.
\end{enumerate}
Only the first implication requires proof. On fresh independent coordinates put
\[
 G_N=N^{-1/2}\sum_{j=1}^N H(X_j).
\]
Then $q+\sqrt vG_N$ is still a finite $H_{2,1}$ polynomial, with
variance $\E|q|^2+v$, and converges in law to $q+\sqrt vG$ by the
two-dimensional central limit theorem. Its uniform second-moment
bound gives uniform integrability of the absolute value. Pass to the
limit in assertion 1. The converse takes $v=0$.

The same argument applies to real third chaos. Thus the Gaussian-buffer
PVD statement in the preceding chapter is equivalent to the original
sharp dimension-free complex $L^1$ inequality; it is not a weaker
technical lemma. The visible Gaussian variance-mixture representation
makes it sufficient for remote K, but dependence of the A/B sectors
on the visible coordinate remains outside this reduction.

\subsection{A phase-sensitive Gaussian-factor dual: proved local and tail parts}
\label{sec:shared-factor-local-dual}
This section records useful finite-dimensional reductions. Its compact
global cover is unfinished, so no global Gaussian-factor A/B theorem
is asserted here.
For deterministic $s,\alpha,\beta\in\mathbb C$, $g\ge0$, define
\[
 R_T=(s+\alpha A(T))X+\beta B(T)\overline X+gG,
 \quad V_0=|s|^2+g^2+|\alpha|^2+|\beta|^2,
\]
\[
 D=2\E|\ell(T)+R_T|-V_0/\varepsilon
      +\frac75\Re(s\overline\alpha)+\frac57\Re(\beta s).
\]
Its actual conditional covariance data are
\[
 V(T)=g^2+|s+\alpha A|^2+|\beta|^2B^2,
 \qquad P(T)=2\beta B(s+\alpha A).
\]
If $D\le S+\eta$ were proved for every deterministic parameter,
then conditioning and averaging would prove the same comparison for
every centered latent $s$ independent of $X,G,T$: both supporting
linear terms have mean zero. Maximizing the orientation of $P(T)$
separately would destroy this cancellation and is not legitimate.

\paragraph{Unbounded parameters.}
The supporting quadratic form is at most $4V_0/5$, since
$49/25+25/49=3026/1225<64/25$. Cauchy--Schwarz gives
\[
 D\le2\sqrt{19/25+V_0}-(1/\varepsilon-4/5)V_0<27/20<S
 \quad(V_0\ge10).
\]
The final inequality and monotonicity on this range are exact rational
comparisons in the supplied checker.

\paragraph{Regression Hessian lemma.}
If centered real planar Gaussian vectors $R,D_0$ are jointly Gaussian
and $\operatorname{Cov}(R)\ge\lambda I$, then
\[
 \sup_z\E\frac{|D_0|^2}{|z+R|}
 \le\frac8{3\sqrt\lambda}\E|D_0|^2.
\]
Write $R=C^{1/2}Y$, $D_0=KY+\xi$ by Gaussian regression.
For a unit vector $e$, the pointwise density bound
$(e\cdot y)^2\phi_I(y)\le(8/e)\phi_{2I}(y)$ and
$\sup_z\E|z+\sqrt2Y|^{-1}=\sqrt\pi/2$ give
$4\sqrt\pi/e<8/3$ for every row of $K$.
The independent $\xi$ term uses
$\sup_z\E|z+Y|^{-1}=\sqrt{\pi/2}<8/3$.
Scaling by the least covariance eigenvalue proves the claim.
Regularization of the norm consequently yields
\[
 \frac{d^2}{dt^2}2\E|z+R+tD_0|
 \le\frac{16}{3\sqrt\lambda}\E|D_0|^2
\]
whenever the covariance along the segment stays above $\lambda I$.

\paragraph{Uniform local support.}
The exact integral checker proves
\[
 |7/5+2\E[A w_v]|\le1/10,\qquad
 |5/7+2\E[B k_v]|\le1/10,
\]
and $|\E[ABk_v]|\le17/100$, for $7/200\le v\le9/200$.
Here $w_v$ and $k_v$ are respectively the angular mean and its
$\cos(2\theta)$ coefficient of
$\sqrt\pi\exp(-\ell(T)^2\cos^2\theta/v)/(2\sqrt v)$.
It follows that
\[
 \boxed{D\le S+\frac3{500000}}
\]
whenever $7/200\le|s|^2+g^2\le9/200$ and
$|\alpha|,|\beta|\le10^{-4}$, uniformly in phases and including $g=0$.
Indeed the supported first derivative is at most $1/200000$.
On $T\le32$, the covariance along the coupling segment is at least
$vI/8$, so the regression bound gives Taylor remainder
$45(|\alpha|^2+|\beta|^2)\le9/10^7$.
The omitted tail of payoff difference minus first derivative is below
$6\cdot10^{-12}$ by the norm Lipschitz bound and exponential moments.
Their sum is less than $3/500000$; discarding the negative coupling
price only increases the bound.

\subsection{Uniform transport of the finite pure-cubic examples}
The finite examples extend to a profile rectangle without another
high-dimensional calculation. Keep $\varepsilon=1891/2500$, and let
$S_{a,b}$ denote the Gaussian benchmark for $L=Z(a+b|Z|^2)$.
For every fixed residual construction $R$ of the types above,
\[
 \Lambda_{a,b}(R)-S_{a,b}
 \le\Lambda_{a_0,b_0}(R)-S_{a_0,b_0}
     +2\sqrt\pi|a-a_0|+3\sqrt\pi|b-b_0|.
\]
Indeed the two-sign payoff changes by at most $2\E|L-L_0|$;
the supremum defining the benchmark changes by the same bound.
Use $\E|Z|=\sqrt\pi/2$ and $\E|Z|^3=3\sqrt\pi/4$.
There is no need to know the maximizing Gaussian variance.

For the main finite family $N\ge66,M\ge20000$, exact squaring of
the buffer error gives the stronger retained peak deficit
\[
 \frac{9221}{250000}-\frac{8987}{5000\cdot66}
     -\frac{9179}{10^6}=\frac{283}{600000},
 \qquad 2\frac9{50}\sqrt{\frac{13}{20000}}<\frac{9179}{10^6}.
\]
Therefore, throughout
\[
 |a-a_0|,|b-b_0|\le\frac1{25000},
\]
the actual finite-cubic family still satisfies
$\Lambda_{a,b}(R)<S_{a,b}-10^{-4}$.
This follows from $\sqrt\pi<1773/1000$ and
\[
 \frac{283}{600000}-\frac{5\cdot1773}{1000\cdot25000}>10^{-4}.
\]
For the optional B amplitude $3/5$, the peak deficit exceeds
$5171/6600000$; the same argument permits half-width $3/40000$
while retaining $10^{-4}$. The rational comparisons are replayed by
\texttt{check\_transport.py}. This is profile uniformity for these
proved families, not a complete all-residual row envelope.

\subsection{Verification and the remaining endpoint obligation}
The package contains exact rational scalar checks, exact polynomial
coefficient checks, and the signed Gaussian integral checker.
It also bundles the inherited characteristic-function and Gaussian
profile certificates, with separate replay commands and provenance.
The current results are mathematical proofs with external numerical
certificates; they have not been formalized in Lean.

The block criterion is a proved sufficient condition for K, and the
finite-amplitude Wick results cover additional cases. A general
coupled cubic does not necessarily decompose into independent blocks.
Nor do these statements settle arbitrary simultaneous even-sector
interactions or a complete admissible row envelope. Therefore no
larger Grothendieck lower endpoint follows from this package alone.
